% Source coverage: physical PDF pp. 473--478 (printed pp. 450--455),
% beginning at the Section 23.5 heading and ending immediately before
% Section 23.6. Numbered equations: (23.5.1)--(23.5.26).
\subsection{Instantons}
\label{sec:23.5}

As we saw in \hyperref[sec:23.1]{Section 23.1}, the topologically non-trivial solutions of a
pure gauge theory with a simple gauge group $G$ in $d=4$ Euclidean
spacetime dimensions correspond to the elements of the homotopy group
$\pi_3(G)=\mathbb{Z}$. These are four-dimensional field configurations, known as
instantons, that (for reasons discussed in the next section) must be
included along with their fluctuations in path integrals. After Belavin,
Polyakov, Schwarz, and Tyupkin~\hyperref[ch23-ref-1]{[1]} demonstrated the
existence of these solutions, 't Hooft~\hyperref[ch23-ref-23]{[23]} showed
that the inclusion of these configurations in path integrals solved the
$U(1)$ problem outlined in \hyperref[sec:19.10]{Section 19.10}. Here we shall first discuss the
instantons themselves, and then consider their role in path integrals.

According to Eq.~\eqref{eq:23.1.7}, in order for a topologically
non-trivial gauge field to have finite action, the gauge fields must
approach a pure gauge at $r\to\infty$ (where here
$r=\sqrt{x_ix_i}$, with $i$ summed over the values 1, 2, 3, 4):
\begin{equation}
iA_i(x)\longrightarrow g^{-1}(\widehat x)\partial_i g(\widehat x),
\tag{23.5.1}\label{eq:23.5.1}
\end{equation}
where $A_i\equiv t_\alpha A_{\alpha i}$ and $g(\widehat x)$ is a
direction-dependent element of the gauge group $G$. The topological
invariant discussed in the previous section may therefore be written in
terms of the asymptotic behavior of the gauge field
\begin{align}
\mathcal I[g]\equiv{}&
\int d\theta^1\,d\theta^2\,d\theta^3\,\epsilon^{abc}
\operatorname{Tr}\left\{
g^{-1}(\theta)\frac{\partial g(\theta)}{\partial\theta^a}
g^{-1}(\theta)\frac{\partial g(\theta)}{\partial\theta^b}
g^{-1}(\theta)\frac{\partial g(\theta)}{\partial\theta^c}
\right\}
\notag\\
={}&-i\lim_{r\to\infty}r^3
\int d\theta^1\,d\theta^2\,d\theta^3\,\epsilon^{abc}
\frac{\partial\widehat x_i}{\partial\theta^a}
\frac{\partial\widehat x_j}{\partial\theta^b}
\frac{\partial\widehat x_k}{\partial\theta^c}
\operatorname{Tr}\{A_iA_jA_k\},
\tag{23.5.2}\label{eq:23.5.2}
\end{align}
where $\theta^a$ with $a=1,2,3$ are any three parameters used to specify
the direction of the unit four-vector $\widehat x$. This surface integral
can be evaluated using Gauss's theorem. In analogy with the current
\eqref{eq:22.2.29} in Minkowski spacetime, we can define a current in
Euclidean spacetime
\begin{equation}
G_l\equiv\epsilon^E_{lijk}\left[
A_{\gamma i}F_{\gamma jk}
-\frac13C_{\alpha\beta\gamma}
A_{\alpha i}A_{\beta j}A_{\gamma k}
\right],
\tag{23.5.3}\label{eq:23.5.3}
\end{equation}
whose divergence is
\begin{equation}
\partial_lG_l=\frac12\epsilon^E_{ijkl}F_{\alpha ij}F_{\alpha kl}.
\tag{23.5.4}\label{eq:23.5.4}
\end{equation}
(Here $\epsilon^E_{ijkl}$ is the totally antisymmetric tensor with
$\epsilon^E_{1234}=1$.) We are using a representation of the gauge group
with totally antisymmetric structure constants, so that
\begin{equation}
\operatorname{Tr}\{t_\alpha t_\beta\}=\frac12N\delta_{\alpha\beta},
\tag{23.5.5}\label{eq:23.5.5}
\end{equation}
with $N$ a constant that depends on the representation in which we
calculate the trace in Eq.~\eqref{eq:23.5.2}. Hence
Eq.~\eqref{eq:23.5.4} may be written
\begin{equation}
G_l=(2/N)\epsilon^E_{lijk}\operatorname{Tr}\left[
A_iF_{jk}+(2i/3)A_iA_jA_k
\right].
\tag{23.5.6}\label{eq:23.5.6}
\end{equation}
For $r\to\infty$, the field strength $F_{kl}$ vanishes, so
\begin{equation}
G_l\longrightarrow(4i/3N)\epsilon^E_{lijk}
\operatorname{Tr}[A_iA_jA_k].
\tag{23.5.7}\label{eq:23.5.7}
\end{equation}
Hence Eq.~\eqref{eq:23.5.2} gives
\begin{equation}
\mathcal I[g]=-\frac{3N}{4}\int d^4x\,\partial_lG_l
=-\frac{3N}{8}\epsilon^E_{ijkl}
\int d^4x\,F_{\alpha ij}F_{\alpha kl}.
\tag{23.5.8}\label{eq:23.5.8}
\end{equation}
Thus in order to demonstrate the existence of topologically non-trivial
field configurations, we have to show that there are configurations for
which the Chern--Pontryagin density
$\epsilon^E_{ijkl}F_{\alpha ij}F_{\alpha kl}$ has a non-vanishing
integral.

For this purpose, it is very useful to take advantage of what is known as
the Bogomol'nyi inequality.~\hyperref[ch23-ref-10]{[10]} From the fact
that
\[
0\leq\int\left(
F_{\alpha ij}\mp\frac12\epsilon^E_{ijkl}F_{\alpha kl}
\right)^2d^4x
\]
(with squares indicating obvious index contractions), we have
\begin{equation}
S[A]\geq\frac18\left|
\epsilon^E_{ijkl}\int F_{\alpha ij}F_{\alpha kl}\,d^4x
\right|
=|\mathcal I[g]|/3N,
\tag{23.5.9}\label{eq:23.5.9}
\end{equation}
where $S[A]$ is (up to a factor) the Euclidean action
\begin{equation}
S[A]\equiv\frac14\int F_{\alpha ij}F_{\alpha ij}\,d^4x.
\tag{23.5.10}\label{eq:23.5.10}
\end{equation}
The lower bound \eqref{eq:23.5.9} is evidently reached if and only if the
gauge field is self-dual or anti-self-dual, in the sense that
\begin{equation}
F_{\alpha ij}=\pm\frac12\epsilon^E_{ijkl}F_{\alpha kl}.
\tag{23.5.11}\label{eq:23.5.11}
\end{equation}
Hence any solution of the first-order equation \eqref{eq:23.5.11} is a
minimum of $S[A]$ for gauge fields of winding number unity, and hence
also a solution of the second-order Yang--Mills field equation.

Belavin \emph{et al.}~\hyperref[ch23-ref-1]{[1]} found a solution of
Eq.~\eqref{eq:23.5.11} of the form
\begin{equation}
iA_i(x)=\left(\frac{r^2}{r^2+R^2}\right)
g_1^{-1}(\widehat x)\partial_i g_1(\widehat x),
\tag{23.5.12}\label{eq:23.5.12}
\end{equation}
where $R$ is an arbitrary scale factor, and $g_1(\widehat x)$ is an
element of an $SU(2)$ subgroup of the gauge group, with
\begin{equation}
g_1(\widehat x)=\left(\frac{x_4+2i\mathbf x\cdot\mathbf t}{r}\right),
\tag{23.5.13}\label{eq:23.5.13}
\end{equation}
and
\begin{equation}
t_1=\frac12\begin{pmatrix}0&1\\1&0\end{pmatrix},
\qquad
t_2=\frac12\begin{pmatrix}0&-i\\i&0\end{pmatrix},
\qquad
t_3=\frac12\begin{pmatrix}1&0\\0&-1\end{pmatrix}.
\tag{23.5.14}\label{eq:23.5.14}
\end{equation}
It is clear that this solution has the asymptotic behavior
\eqref{eq:23.5.1}, with $g(\widehat x)$ the same as the `identity'
mapping \eqref{eq:23.4.9}, so this solution belongs to the homotopy class
of the identity map, and therefore as we saw in the previous section has
winding number $\nu=1$. From Eqs.~\eqref{eq:23.4.14} and
\eqref{eq:23.5.9} (which is here an equality) we have then
\begin{equation}
S[A]=8\pi^2.
\tag{23.5.15}\label{eq:23.5.15}
\end{equation}
From Eqs.~\eqref{eq:23.5.10} and \eqref{eq:23.5.11} (with positive
sign), we also have
\begin{equation}
\epsilon^E_{ijkl}\int F_{\alpha ij}F_{\alpha kl}\,d^4x=64\pi^2.
\tag{23.5.16}\label{eq:23.5.16}
\end{equation}
This solution is not unique, because it can be translated or subjected
to a gauge transformation, but aside from these degrees of freedom,
there are no other solutions of the field equations with winding number
unity.~\hyperref[ch23-ref-24]{[24]}

Because we have found field configurations with $\nu=1$, we know that
there are also field configurations with any integer $\nu$. For instance,
solutions with $\nu$ a positive integer $\mathcal N$ can be constructed
by superimposing $\mathcal N$ solutions with $\nu=1$ with centers so far
apart that at these distances the non-linearities of the field equations
become unimportant. A solution with $\nu=-1$ negative can be found by
replacing $g_1$ in Eq.~\eqref{eq:23.5.12} with $g_1^{-1}$, and solutions
with $\nu$ a negative integer $-\mathcal N$ can be found by
superimposing $\mathcal N$ of these at large separations. For general
winding number, Eqs.~\eqref{eq:23.5.15} and \eqref{eq:23.5.16} become
\begin{equation}
S[A]=8\pi^2|\nu|,
\tag{23.5.17}\label{eq:23.5.17}
\end{equation}
\begin{equation}
\epsilon^E_{ijkl}\int F_{\alpha ij}F_{\alpha kl}\,d^4x=64\pi^2\nu.
\tag{23.5.18}\label{eq:23.5.18}
\end{equation}

These results are for a gauge field normalized as in
Eqs.~\eqref{eq:23.5.11}--\eqref{eq:23.5.13}. With this normalization,
the action $I[A]$ is not $-S[A]$, but
\begin{equation}
I[A]=-S[A]/g^2=-8\pi^2|\nu|/g^2,
\tag{23.5.19}\label{eq:23.5.19}
\end{equation}
where $g$ is the conventional coupling constant. If we had used our usual
convention of including a factor $g$ in the generators and structure
constant, then the action $I[A]$ would be the same as $-S[A]$, but with
$A_{\alpha\mu}$ and $F_{\alpha\mu\nu}$ carrying a factor $1/g$ we would
have $S[A]=8\pi^2/g^2$ instead of Eq.~\eqref{eq:23.5.14}, so in this
case again the action would be $-8\pi^2/g^2$, but now in place of
Eq.~\eqref{eq:23.5.14} we would have
\begin{equation}
\epsilon^E_{ijkl}\int F_{\alpha ij}F_{\alpha kl}\,d^4x
=64\pi^2\nu/g^2.
\tag{23.5.20}\label{eq:23.5.20}
\end{equation}

We shall see in the next section that in path integrals we must sum the
effects of instantons of all winding numbers. The contribution of
configurations of winding number $\nu\ne0$ to Euclidean path integrals is
suppressed by a factor
$\exp(I[A])=\exp(-8|\nu|\pi^2/g^2)$. In \hyperref[sec:23.7]{Section 23.7} we will see that
the coefficient of this exponential is a negative power $-n$ of $g$: in
quantum chromodynamics, $n=12$. The function
$g^{-n}\exp(-8|\nu|\pi^2/g^2)$ and all its derivatives with respect to
$g$ vanish at $g=0$, so such contributions are non-perturbative---they
will never be encountered in any order of perturbation theory.

This does not necessarily mean that these contributions are small. As we
saw in Chapter 18, in quantum chromodynamics the coupling $g$ is not a
fixed dimensionless parameter, but a function of a sliding energy scale,
and becomes large at low energies. The effective energy scale to use in
the coupling in Eq.~\eqref{eq:23.5.19} is determined by quantum
fluctuations, the subject of \hyperref[sec:23.7]{Section 23.7}, but on dimensional grounds it
cannot be very different from $1/R$, where $R$ is the instanton size in
Eq.~\eqref{eq:23.5.12}. The instanton size is not fixed, but must be
integrated over, with some weight function that depends on the process
under consideration. In quantum chromodynamics the running coupling
constant $g_\mu$ is given for large $\mu$ by Eq.~\eqref{eq:18.7.7} as
\[
g_\mu^2=\frac{8\pi^2}{\beta_0\ln(\mu/\Lambda)},
\]
where $\beta_0=11-2n_f/3$ and $\Lambda\approx250\ \mathrm{MeV}$ is the
quantum chromodynamic scale factor. The factor
$\exp(-8\pi^2/g_{1/R}^2)$ is therefore, for small instantons,
\[
\exp\left(-8\pi^2/g_{1/R}^2\right)=(R\Lambda)^{\beta_0}.
\]
We cannot calculate this factor for large instantons, with
$R\Lambda\gg1$, but it is clear that in this case there is no suppression
of instanton effects.

Despite initial hopes, the discovery of instantons has not led to much
improvement in our ability to do quantitative calculations in quantum
chromodynamics. On the other hand, as we shall now see, it has produced
spectacular \emph{qualitative} changes in our understanding of quantum
chromodynamics and other gauge theories.

The mere fact that there are solutions of the field equations for which
the integral \eqref{eq:23.5.15} does not vanish is enough to provide a
solution of the $U(1)$ problem discussed in \hyperref[sec:19.10]{Section 19.10}. Under
\emph{global} $U(1)$ transformations
$\psi\to\exp(i\gamma_5\alpha)\psi$, the measure for integration over
quark fields undergoes a change given by Eq.~\eqref{eq:22.2.10}:
\begin{equation}
[d\psi][d\overline\psi]\longrightarrow
\exp\left\{i\alpha\int\mathcal A(x)\,d^4x\right\}
[d\psi][d\overline\psi],
\tag{23.5.21}\label{eq:23.5.21}
\end{equation}
where (with the matrix $t$ there equal to unity) the anomaly function
$\mathcal A(x)$ is given by Eq.~\eqref{eq:22.2.45}:
\begin{equation}
\mathcal A(x)=\frac{1}{16\pi^2}\epsilon^E_{ijkl}
F_{ij\alpha}F_{kl\beta}\operatorname{tr}t_\alpha t_\beta.
\tag{23.5.22}\label{eq:23.5.22}
\end{equation}
The existence of this anomaly would not in itself solve the $U(1)$
problem, because the quantity \eqref{eq:23.5.22} is a total derivative,
and so would have a zero integral for a non-singular gauge field that
vanishes sufficiently rapidly at infinity. The instanton solution
vanishes only as $1/r$, and yields a non-vanishing value for this
integral, given by Eqs.~\eqref{eq:23.5.22}, \eqref{eq:23.5.18}, and
\eqref{eq:23.5.5} as
\begin{equation}
\int d^4x\,\mathcal A(x)=2N\nu,
\tag{23.5.23}\label{eq:23.5.23}
\end{equation}
showing that the anomaly \emph{does} violate the $U(1)$ chiral symmetry.

As we saw in \hyperref[sec:22.4]{Section 22.4}, the currents of baryon and lepton number also
contain anomalies due to the interaction of quarks and leptons with the
$SU(2)\times U(1)$ gauge fields of the standard model. Instanton
configurations of the $SU(2)$ gauge field therefore produce violations
of baryon and lepton conservation.~\hyperref[ch23-ref-25]{[25]} As noted
in \hyperref[sec:22.4]{Section 22.4}, there are various currents whose conservation is not
violated by anomalies or anything else, such as baryon number minus
lepton number, and the differences of electron, muon, and tau lepton
numbers, so these will be conserved in any baryon- and
lepton-non-conserving process. For instance, the decay of a proton or
deuteron is forbidden, but the decay
$\mathrm{He}^3\to e^++\mu^++\overline\nu_\tau$ is allowed. The
amplitudes for these effects are suppressed by the same factor
$\exp(-8|\nu|\pi^2/g^2)$ as before, but now with $g$ the $SU(2)$
coupling constant $e/\sin\theta$, evaluated not at a sliding scale but at
the natural scale for electroweak processes, roughly of order $m_Z$.
Taking $e^2/4\pi=1/129$ (see \hyperref[sec:18.2]{Section 18.2}) and
$\sin^2\theta=0.23$, the suppression factor for $|\nu|=1$ is
$\exp(-373)$. $\mathrm{He}^3$ decay into three antileptons is not likely
to be observed.

\begin{center}
$* \quad * \quad *$
\end{center}

There is another approach~\hyperref[ch23-ref-26]{[26]} that illuminates
some aspects of instantons. In temporal gauge, with
$A_{\alpha4}(\mathbf x,x_4)=0$, the gauge field for
$x_4\to\pm\infty$ is expected to approach time-independent pure gauges
for $x_4\to\pm\infty$:
\begin{equation}
i\mathbf A(\mathbf x,x_4)\equiv
it_\alpha\mathbf A_\alpha(\mathbf x,x_4)
\longrightarrow
g_\pm^{-1}(\mathbf x)\boldsymbol{\nabla}g_\pm(\mathbf x),
\tag{23.5.24}\label{eq:23.5.24}
\end{equation}
where $g_\pm(\mathbf x)$ are group elements in the representation
generated by $t_\alpha$. Assuming that
$\mathbf A_\alpha(\mathbf x,x_4)$ vanishes for $\mathbf x\to\infty$, the
group elements $g_\pm(\mathbf x)$ must approach constants $g_\pm$ for
$\mathbf x\to\infty$, so the three-spaces at $x_4\to\infty$ may be
regarded as three-spheres, with the point at infinity regarded as an
ordinary point. Following the same argument that led to
Eq.~\eqref{eq:23.4.13}, we then have
\begin{equation}
\epsilon_{ijk}\int d^3x\,
\operatorname{Tr}\left\{
g_\pm^{-1}(\mathbf x)\partial_i g_\pm(\mathbf x)
g_\pm^{-1}(\mathbf x)\partial_j g_\pm(\mathbf x)
g_\pm^{-1}(\mathbf x)\partial_k g_\pm(\mathbf x)
\right\}
=24\pi^2n_\pm,
\tag{23.5.25}\label{eq:23.5.25}
\end{equation}
where $n_\pm$ are integers. The integral over the boundary of four-space
in Eq.~\eqref{eq:23.5.2} may be regarded as the integral over the `plane'
$x_4=+\infty$ minus the integral over the `plane' $x_4=-\infty$, so from
\eqref{eq:23.5.8} and \eqref{eq:23.5.18} (with $N=1$) we have
\begin{equation}
\nu=n_+-n_-.
\tag{23.5.26}\label{eq:23.5.26}
\end{equation}
The exponential factor $\exp(-8|\nu|\pi^2/g^2)$ from
Eq.~\eqref{eq:23.5.19} may therefore be regarded as the amplitude for a
transition from a configuration with spatial winding number $n_-$ at
$x_4\to-\infty$ to one of spatial winding number $n_+$ at
$x_4\to+\infty$. The exponential form of this factor for $\nu\ne0$
reflects the fact that this is a \emph{tunneling} process; no continuous
sequence of pure gauge fields can take us from a configuration of one
spatial winding number to a configuration with a different spatial
winding number.

The interpretation of the factor $\exp(-8|\nu|\pi^2/g^2)$ as a tunneling
amplitude suggests that baryon- and lepton-non-conserving processes may
proceed rapidly at temperatures above about 1 TeV, where instead of
having to tunnel through the barrier, thermal fluctuations can take the
vacuum \emph{over} the barrier.~\hyperref[ch23-ref-27]{[27]} This process
may be of cosmological importance, but it is still subject to the
selection rules mentioned above: thermal fluctuations will not change
the density of baryon number minus lepton number, or the differences of
the densities of the three varieties of lepton number.
